\documentclass[11pt]{article}

\usepackage[margin=1in]{geometry}
\usepackage{amsmath,amssymb,amsthm,mathtools}
\usepackage{enumitem}
\usepackage{booktabs}
\usepackage{tabularx}
\usepackage{microtype}
\usepackage[hidelinks]{hyperref}
\usepackage{xurl}

\newtheorem{definition}{Definition}
\newtheorem{lemma}{Lemma}
\newtheorem{remark}{Remark}

\title{Public Request Response Packet Carryforward for Source-Only Successor Releases:\\[0.5ex]
\large Reuse Rules for Exact Notice-Keyed Handoff Slices under Release Evolution}
\author{}
\date{Unified archive note v0.08 (2026-03-21; archive v476)}

\begin{document}
\maketitle

\begin{abstract}
The source-only branch can now say what a terse paper promises on request, whether that paper-level promise survives a successor cut, what exact visible notice readers saw, which canonical notice family and selector justified that notice, which already-maintained owner is the smallest sufficient answer for one concrete follow-up about that notice, and what exact bounded packet should actually travel once that owner has been chosen.
One cross-release question still remained.
When a successor release ships, may that exact packet be reused as is, refreshed in place, or must it reopen?
This note gives that successor rule one home.
It defines compact \emph{public request response packet carryforward profiles}: successor-release maintenance objects keyed to one already-emitted source-only response packet and one successor carryforward context, classifying whether that packet survives unchanged, survives with refreshable release-bound fields, or reopens.
\end{abstract}

\paragraph{Non-redundancy note.}
Synthesis~60 owns whether the paper-level source-only promise itself survives, Synthesis~64 owns which already-maintained owner is the smallest sufficient answer for one follow-up class, Synthesis~65 owns the exact bounded packet cut that should travel once that owner has been chosen, Synthesis~73 owns the request-side terminal stop later terse papers should usually cite, Synthesis~74--75 own the paired terminal / cutover discipline, and Synthesis~17 owns the concrete worked instantiation.\cite{series:publicrequestcarryforward,series:publicrequestresponsemenus,series:publicrequestresponsepackets,series:publicrequestresponsepacketrefreshresponsepacketclosureverdicts,series:pairedterminalcitationdiscipline,series:pairedterminalcutoverrule,series:workedexample}
This note owns only the successor-release reuse discipline for one already-cut response packet. Repeated refreshable-field-family labels may recur across many such profiles, but only as imported packet-maintenance vocabulary; the reuse verdict itself remains packet-scoped, and later terse papers should name it directly only when packet reuse itself is the live issue.

\paragraph{Scope.}
This is not a new carryforward verdict for the paper-level promise, not a new response-menu choice, and not a new bounded packet cut.
It is a thin successor-release bridge answering: after a concrete source-only response packet has already been fixed, may that exact handoff slice be reused as is, refreshed in place, or must it reopen under the successor cut? In the worked route/spine bridge it is also the first reopened owner when a visible-refresh closure fails because packet reuse itself has flipped back to \texttt{reopen}.

\paragraph{Exports.}
This note exports (i) public request response packet carryforward profiles, (ii) a preservation rule, (iii) a refresh-in-place rule, (iv) a shared-vocabulary rule, (v) a reopen rule, and (vi) a mirrored packet-maintenance-posture rule saying that later terse papers may cite one compact maintained profile when the exact successor reuse status of an exact source-only handoff slice is itself the live issue. If they instead need only the stable request-side stopping point or the stop-versus-widen judgment for that branch, they should cite Synthesis~73--75 rather than restitch this packet-reuse layer locally.

\paragraph{Later-terse-paper citation posture.}
Later terse papers should cite Synthesis~66 only when the live issue is the \emph{successor reuse verdict itself} for one already-cut source-only response packet: whether that exact predecessor packet remains \texttt{preserved}, remains adequate but only as \texttt{refresh\_in\_place}, or must be marked \texttt{reopened}. They should stop at Synthesis~60 when the live issue is the paper-level source-only carryforward envelope, stop at Synthesis~64 when the live issue is the smallest-sufficient follow-up owner, stop at Synthesis~65 when the live issue is the exact predecessor packet cut, stop at Synthesis~67 when the live issue is which packet-local bindings refreshed, stop at Synthesis~68--71 when the live issue is the emitted visible refresh sentence, its canonical wrapper family, the family-choice selector, or the smallest-sufficient successor-facing handoff, widen to Synthesis~72 when the live issue is the exact successor-facing packet cut after that reuse verdict is fixed, widen to Synthesis~73--75 when the remaining issue is request-side terminal stop, paired terminal reuse, or fail-closed cutover, and stop at Synthesis~17 when the concrete maintained worked route is the live issue instead.\cite{series:publicrequestresponsepacketrefreshnotices,series:publicrequestresponsepacketrefreshnoticenormalforms,series:publicrequestresponsepacketrefreshnoticeselections,series:publicrequestresponsepacketrefreshresponsemenus,series:publicrequestresponsepacketrefreshresponsepackets}

\section{Why packet carryforward deserves its own note}
A response packet answers the within-release question ``what exact bounded slice should travel right now?''
A successor release adds a different question.
Suppose an older terse paper remains published and readers still ask concrete follow-up questions about the visible source-only notice in that paper.
May the previously certified packet still be handed over, perhaps after refreshing release-bound digests or basis companions, or has the successor cut changed enough that the packet must be rebuilt or reopened?
Without one maintained carryforward object, later notes either restate packet reuse rules in prose or silently assume that a preserved paper-level status token automatically means the whole predecessor packet can be reused unchanged.
Those are different claims.

\begin{definition}[Public request response packet carryforward profile]
Fix one predecessor public request response packet $\Pi$, one successor public request carryforward envelope $K$, and the successor-release continuity context from the maintained release spine.
A \emph{public request response packet carryforward profile} is a tuple
\[
\Gamma(\Pi,K)=\bigl(\Pi,K,\rho,\Phi,\Delta\bigr),
\]
where $\rho$ is one reuse token with allowed values
\[
\texttt{preserved},\qquad \texttt{refresh\_in\_place},\qquad \texttt{reopened},
\]
and $\Phi$ and $\Delta$ record the refreshable-field-family list and changed-context-class list for that same predecessor packet key.
The profile does not replace packet ownership.
It only states how that already-owned packet behaves when the successor release becomes the current maintenance context.
\end{definition}

\begin{definition}[Preservation rule]
A packet carryforward profile satisfies the \emph{preservation rule} if it emits \texttt{preserved} exactly when the imported successor context leaves the selected owner, bounded companion set, and packet meaning unchanged and no release-bound field in the packet requires refresh.
So preservation is stronger than paper-level status carryforward: it says the exact bounded handoff slice itself remains ready to ship without modification.
\end{definition}

\begin{definition}[Refresh-in-place rule]
A packet carryforward profile satisfies the \emph{refresh-in-place rule} if it emits \texttt{refresh\_in\_place} exactly when the selected owner and bounded slice remain sufficient under the successor context, but some release-bound fields may need refreshing in place.
Permitted refreshes are limited to the carried basis companions, release/digest bindings, or packet excerpts whose owning narrower object already allows successor refresh without changing the follow-up-class owner or widening the packet.
So refresh-in-place is a reuse verdict, not a license to silently reroute the handoff.
\end{definition}

\begin{definition}[Shared-vocabulary rule]
A packet carryforward profile satisfies the \emph{shared-vocabulary rule} if any repeated refreshable-field-family label in $\Phi$ denotes only imported packet-maintenance vocabulary that may recur across many profiles, while the reuse token $\rho$ and adequacy claim still belong to the concrete predecessor packet named in that same profile.
So recurring family labels do not silently turn a packet-scoped reuse verdict into a family-scoped one.
\end{definition}

\begin{definition}[Reopen rule]
A packet carryforward profile satisfies the \emph{reopen rule} if it emits \texttt{reopened} whenever the successor carryforward context changes the selected owner, changes the required companion set, or reopens a promised family so that the predecessor packet is no longer the smallest sufficient bounded handoff.
So reopening belongs to packet adequacy under successor maintenance, not merely to the paper-level source-only token.
\end{definition}

\begin{definition}[Mirrored packet-maintenance-posture rule]
A packet carryforward profile satisfies the \emph{mirrored packet-maintenance-posture rule} if it records the current packet reuse posture once, and every later delta / refresh-notice / successor-packet / route-side owner mirrors that posture exactly rather than reconstructing it ad hoc from the reuse token alone.
So the concrete packet-scoped reuse verdict stays visible even when later notes compress it into narrower maintenance objects.
\end{definition}

\begin{lemma}[Packet carryforward profiles stabilize successor handling of exact source-only handoff slices without replacing packet owners]
Assume every public request response packet carryforward profile satisfies the preservation, refresh-in-place, shared-vocabulary, reopen, and mirrored packet-maintenance-posture rules.
Then a later terse paper may cite one maintained profile to explain whether a previously shipped exact source-only response packet still travels unchanged, may be refreshed in place, or must reopen after a successor release without re-owning the packet cut, the response-menu choice, or the underlying paper-level carryforward verdict.
\end{lemma}

\begin{proof}
The predecessor response packet still owns the exact bounded handoff slice.
The paper-level carryforward envelope still owns whether the underlying source-only promise remained complete, advanced, or reopened.
The carryforward profile reads those maintained owners and adds only one new statement: the successor-release reuse class of that already-owned packet.
Because the refresh-in-place rule limits successor edits to release-bound fields that the imported narrower owners already permit, the shared-vocabulary rule keeps repeated family labels from becoming family-scoped verdict owners, the reopen rule fires whenever owner choice or required companions change, and the mirrored packet-maintenance-posture rule keeps later maintenance objects from quietly restating the verdict differently. The profile therefore cannot silently replace the packet or widen it while pretending only to maintain reuse state.
So later papers get one stable successor-handling answer without collapsing packet ownership into the release spine.\qedhere
\end{proof}

\begin{table}[t]
\centering
\footnotesize
\setlength{\tabcolsep}{2pt}
\begin{tabularx}{\linewidth}{@{}>{\raggedright\arraybackslash}p{0.18\linewidth}>{\raggedright\arraybackslash}p{0.20\linewidth}>{\raggedright\arraybackslash}X@{}}
\toprule
Reuse token & Meaning & Consequence \\
\midrule
\texttt{preserved} & Exact bounded packet still ships unchanged under successor maintenance & Reuse the packet as-is. \\
\texttt{refresh\_in\_place} & Same owner and bounded slice remain sufficient, but release-bound fields must be refreshed & Refresh only the listed field families, then reissue the same packet. \\
\texttt{reopened} & Successor context changes owner choice, required companions, or family completion enough that the old packet is no longer sufficient & Rebuild or reroute from the narrower owners instead of reusing the predecessor packet. \\
\bottomrule
\end{tabularx}
\caption{Packet carryforward profiles classify the successor-release reuse status of an already-owned source-only response packet.}
\label{tab:publicrequestresponsepacketcarryforward}
\end{table}

\paragraph{A minimal public packet-carryforward card.}
\begin{table}[t]
\centering
\footnotesize
\begin{tabularx}{\linewidth}{@{}l X@{}}
\toprule
Card field & Worked floor \\
\midrule
Profile keys & Representative worked keys: \nolinkurl{worked_example_receipt_interlock_within_release_public_request_response_menu__basis_token_check__carryforward} and \nolinkurl{worked_example_receipt_interlock_successor_public_request_response_menu__request_scope_check__carryforward}; all twelve worked entries share the same carryforward-token family and differ only by imported predecessor packet / selected owner / notice stem \\
Imported predecessor packet floor & The worked profiles stay attached to already-cut predecessor packets such as \nolinkurl{worked_example_receipt_interlock_within_release_public_request_response_menu__basis_token_check} and \nolinkurl{worked_example_receipt_interlock_successor_public_request_response_menu__request_scope_check}, and keep the selected owner visible through keys like \nolinkurl{worked_example_receipt_interlock_public_request_status_envelope} or \nolinkurl{worked_example_receipt_interlock_public_request_contract} \\
Successor context floor & Every worked profile imports the same successor context \nolinkurl{worked_example_receipt_interlock_public_request_carryforward_envelope} with \nolinkurl{continuity_decision = same_certified_interface_replayed_surfaces} \\
Reuse verdict floor & All twelve worked entries emit \nolinkurl{response_packet_carryforward_token = refresh_in_place}; their mirrored packet-maintenance posture stays \nolinkurl{current_status = refresh_in_place} and \nolinkurl{reissue_status = refresh_then_reissue} \\
Refreshable-family floor & Every worked entry exposes exactly three refreshable field families---\nolinkurl{basis_companions}, \nolinkurl{selected_entry_excerpt}, and \nolinkurl{digest_bindings}---under \nolinkurl{refreshable_field_family_mode = shared_vocabulary_packet_scoped_verdict} \\
Changed-context floor & The canned successor context changes only \nolinkurl{successor_release_binding} and \nolinkurl{successor_digest_binding}; no worked profile changes the selected owner, request class, or bounded packet membership \\
Escalation pointer & Stop at Synthesis~65 when the live issue is still the exact predecessor packet cut; stop at Synthesis~67 when the live issue is the exact refreshed field list; widen to Synthesis~68--72 only when the visible refresh wrapper or successor-facing packet cut is already the live issue; widen to Synthesis~73--75 only for request-side terminal / paired cutover ownership \\
\bottomrule
\end{tabularx}
\caption{A minimal public packet-carryforward card. This is the smallest public answer later terse papers can quote directly when the live issue is whether one already-cut source-only response packet is preserved, refreshed in place, or reopened across one concrete successor release rather than the narrower predecessor packet cut, the field-level delta ledger, or the wider successor refresh wrapper tail.}
\label{tab:minimal-public-packet-carryforward-card}
\end{table}

This card is intentionally non-decorative: it pins the actual worked carryforward-profile keys, the imported predecessor packet and successor-envelope floors, the exact refreshable family vocabulary, and the mirrored packet-maintenance posture without silently re-owning the predecessor packet cut, the packet-local delta ledger, or the successor-facing refresh wrapper family that neighboring notes already own.\cite{series:workedexample}


\paragraph{A forced public packet-carryforward matrix.}
\begin{table}[t]
\centering
\footnotesize
\begin{tabularx}{0.97\linewidth}{@{}p{0.31\linewidth}X X@{}}
\toprule
Worked carryforward profile & Still the same carryforward profile while & Shortcut goes stale when \\
\midrule
\nolinkurl{worked_example_receipt_interlock_within_release_public_request_response_menu__basis_token_check__carryforward} & Same predecessor packet \nolinkurl{worked_example_receipt_interlock_within_release_public_request_response_menu__basis_token_check}, same imported notice lineage \nolinkurl{worked_example_receipt_interlock_public_request_status_notice}, same successor envelope \nolinkurl{worked_example_receipt_interlock_public_request_carryforward_envelope} with \nolinkurl{continuity_decision = same_certified_interface_replayed_surfaces}, same \nolinkurl{response_packet_carryforward_token = refresh_in_place}, same mirrored posture \nolinkurl{current_status = refresh_in_place} / \nolinkurl{reissue_status = refresh_then_reissue}, and the same three refreshable field families plus the same two changed-context classes. & Any drift in predecessor packet key, selected owner, imported notice lineage, successor-envelope continuity decision, reuse token, mirrored packet-maintenance posture, refreshable-family set, or changed-context-class set means the same basis-token carryforward profile is no longer forced and Synthesis~66 must be reopened. \\
\nolinkurl{worked_example_receipt_interlock_successor_public_request_response_menu__request_scope_check__carryforward} & Same predecessor packet \nolinkurl{worked_example_receipt_interlock_successor_public_request_response_menu__request_scope_check}, same selected owner \nolinkurl{worked_example_receipt_interlock_public_request_contract}, same imported notice lineage \nolinkurl{worked_example_receipt_interlock_public_request_carryforward_notice}, same successor envelope / continuity decision, same \nolinkurl{refreshable_field_families = [basis_companions, selected_entry_excerpt, digest_bindings]}, and the same mirrored packet-maintenance posture. & Any drift in predecessor packet membership, selected owner excerpt, imported notice lineage, successor-envelope floor, mirrored reuse posture, or declared refreshable-family list means the same request-scope carryforward profile is no longer forced; the note must reopen Synthesis~66 instead of inferring packet reuse from the older packet cut alone. \\
\bottomrule
\end{tabularx}
\caption{A forced public packet-carryforward matrix. Once an exact predecessor packet is fixed, the same successor reuse profile remains mechanically forced only while the imported predecessor packet, successor-envelope floor, mirrored packet-maintenance posture, and refreshable-family vocabulary remain unchanged.}
\label{tab:forced-public-packet-carryforward-matrix}
\end{table}

This matrix makes the packet-carryforward stop rule explicit: one predecessor packet key is not enough by itself. A later terse paper gets the same successor reuse profile only while the predecessor packet lineage, imported notice lineage, successor-envelope continuity decision, mirrored reuse posture, and refreshable-family vocabulary still agree. Once any of those drift, the carryforward shortcut has gone stale and Synthesis~66 must be reopened.\cite{series:workedexample}

\paragraph{One tiny packet-carryforward drill.}
For \nolinkurl{worked_example_receipt_interlock_successor_public_request_response_menu__request_scope_check__carryforward}, the predecessor packet is still \nolinkurl{worked_example_receipt_interlock_successor_public_request_response_menu__request_scope_check}, the selected owner is still \nolinkurl{worked_example_receipt_interlock_public_request_contract}, and the successor context still says \nolinkurl{same_certified_interface_replayed_surfaces}. The only worked differences lie in the release-bound packet families \nolinkurl{basis_companions}, \nolinkurl{selected_entry_excerpt}, and \nolinkurl{digest_bindings}. So the honest successor verdict is \nolinkurl{refresh_in_place}: not \nolinkurl{preserved}, because release-bound packet fields still need rebinding, and not \nolinkurl{reopened}, because neither the selected owner nor the bounded packet membership changed. Later terse papers can therefore quote the carryforward profile directly only as long as the same predecessor packet, same selected owner, and same changed-context floor remain fixed.\cite{series:workedexample}

\section{Worked example pointer}
The maintained worked bundle now ships \nolinkurl{example_public_request_response_packet_carryforward_profiles.json}.\cite{series:workedexample}
It contains twelve carryforward profiles: one for each of the six shared follow-up classes under the worked within-release menu and one for each matching class under the worked successor-carryforward menu.
In the maintained canned successor comparison, every profile records the same packet-scoped successor verdict and the same packet-maintenance posture:
\begin{quote}\small
\nolinkurl{response_packet_carryforward_token=refresh_in_place}.\par
\nolinkurl{continuity_decision=same_certified_interface_replayed_surfaces}.\par
Changed-context classes: \nolinkurl{successor_release_binding}, \nolinkurl{successor_digest_binding}.\par
Refreshable field families: \nolinkurl{basis_companions}, \nolinkurl{selected_entry_excerpt}, \nolinkurl{digest_bindings}.\par
Maintenance posture: \nolinkurl{current_status=refresh_in_place}, \nolinkurl{reissue_status=refresh_then_reissue}.\par
Field-family mode: \nolinkurl{refreshable_field_family_mode=shared_vocabulary_packet_scoped_verdict}.
\end{quote}
Repeated field-family labels therefore recur across the twelve profiles only as imported maintenance vocabulary, while the reuse verdict itself remains attached to each concrete predecessor packet key.
Packet-reuse challenges should therefore stop first at the exact carryforward profile for the visible notice/request-class pair, then widen only to the exact predecessor response packet and imported paper-level carryforward envelope if a referee needs to see why that reuse verdict is justified; only later do they widen to Synthesis~67 when the question becomes which packet-local bindings refreshed and which stayed fixed under reissue.\cite{series:publicrequestresponsepacketdeltas}
But later terse papers that are not disputing packet reuse itself should usually cite the request-side terminal form and paired cutover rule from Synthesis~73--75 rather than narrating this carryforward layer in local prose. The maintained worked route catalog mirrors this directly by importing both the exact packet-carryforward-profile key and a packet-reuse-status tag for each shipped notice/follow-up pair, so an audit-focused downstream note can still recover the visible refresh branch's reuse verdict from the same maintained route object that already names the checked public sentence and bounded omitted-support packet.

\begin{thebibliography}{99}\small

\bibitem{series:publicrequestcarryforward}
Anonymous.
\newblock Public Request Carryforward for Source-Only Successor Releases: When Paper-Level Request Status Persists, Advances, or Reopens across Release Evolution.
\newblock Series synthesis note (Synthesis~60), 2026. (This archive.)

\bibitem{series:publicrequestresponsemenus}
Anonymous.
\newblock Public Request Response Menus for Source-Only Releases: Smallest-Sufficient Follow-Up Owners for Visible Notices.
\newblock Series synthesis note (Synthesis~64), 2026. (This archive.)

\bibitem{series:publicrequestresponsepackets}
Anonymous.
\newblock Public Request Response Packets for Source-Only Releases: Exact Bounded Handoff Slices after Notice-Keyed Follow-Up Choice.
\newblock Series synthesis note (Synthesis~65), 2026. (This archive.)

\bibitem{series:publicrequestresponsepacketrefreshnotices}
Anonymous.
\newblock Public Request Response Packet Refresh Notices for Source-Only Successor Releases: Visible Successor Reissue Sentences after Reused Notice-Keyed Handoff Slices.
\newblock Series synthesis note (Synthesis~68), 2026. (This archive.)

\bibitem{series:publicrequestresponsepacketrefreshnoticenormalforms}
Anonymous.
\newblock Public Request Response Packet Refresh Notice Normal Forms for Source-Only Successor Releases: Canonical Visible Wrapper Families for Refreshed Handoff Slices.
\newblock Series synthesis note (Synthesis~69), 2026. (This archive.)

\bibitem{series:publicrequestresponsepacketrefreshnoticeselections}
Anonymous.
\newblock Public Request Response Packet Refresh Notice Selections for Source-Only Successor Releases: Family-Choice Selectors for Visible Reissue Wrappers.
\newblock Series synthesis note (Synthesis~70), 2026. (This archive.)

\bibitem{series:publicrequestresponsepacketrefreshresponsemenus}
Anonymous.
\newblock Public Request Response Packet Refresh Response Menus for Source-Only Successor Releases: Smallest-Sufficient Successor Handoffs after Visible Reissue.
\newblock Series synthesis note (Synthesis~71), 2026. (This archive.)

\bibitem{series:publicrequestresponsepacketrefreshresponsepackets}
Anonymous.
\newblock Public Request Response Packet Refresh Response Packets for Source-Only Successor Releases: Exact Successor-Facing Handoff Slices after Visible Reissue.
\newblock Series synthesis note (Synthesis~72), 2026. (This archive.)

\bibitem{series:publicrequestresponsepacketdeltas}
Anonymous.
\newblock Public Request Response Packet Delta Ledgers for Source-Only Successor Releases: Exact Refreshed Handoff Bindings for Reused Notice-Keyed Packet Slices.
\newblock Series synthesis note (Synthesis~67), 2026. (This archive.)

\bibitem{series:publicrequestresponsepacketrefreshresponsepacketclosureverdicts}
Anonymous.
\newblock Public Request Response Packet Refresh Response Packet Closure Verdicts for Source-Only Successor Releases: Capstone Certificates for Exact Successor-Facing Handoff Slices.
\newblock Series synthesis note (Synthesis~73), 2026. (This archive.)

\bibitem{series:pairedterminalcitationdiscipline}
Anonymous.
\newblock Paired Terminal Citation Discipline for Stable Review Spines: One Answer-Side Stop and One Request-Side Capstone for Terse Papers.
\newblock Series synthesis note (Synthesis~74), 2026. (This archive.)

\bibitem{series:pairedterminalcutoverrule}
Anonymous.
\newblock Paired Terminal Cutover Rules for Stable Review Spines: When Terse Papers Stop, Widen, or Reopen.
\newblock Series synthesis note (Synthesis~75), 2026. (This archive.)

\bibitem{series:workedexample}
Anonymous.
\newblock Worked Example: Receipt Interlock for an Anonymous-DHT Reader-Privacy Claim.
\newblock Series synthesis note (Synthesis~17), 2026. (This archive.)

\end{thebibliography}
\end{document}
